% Part: normal-modal-logic
% Chapter: filtrations
% Section: filtrations-def

\documentclass[../../../include/open-logic-section]{subfiles}

\begin{document}

\olfileid[pt]{nml}{fil}{fil}

\olsection{Filtrações}

Em vez de definir ``a'' filtração de $\mModel{M}$ por $\Gamma$,
definimos quando um modelo~$\mModel{M^*}$ conta como filtração
de~$\mModel{M}$. Todas as filtrações têm o mesmo conjunto de
mundos~$W^*$ e a mesma valoração~$V^*$. Contudo, filtrações diferentes
podem ter relações de acessibilidade~$R^*$ diferentes. Para contar
como filtração, $R^*$ deve satisfazer várias condições. Essas condições
são exatamente o que exigiremos para demonstrar o resultado principal,
a saber, que $\mSat{M}{!A}[w]$ se e somente se
$\mSat{M^*}{!A}[{[w]}]$, desde que $!A \in \Gamma$.

\begin{defn}\ollabel{defn:filtration}
  Seja $\Gamma$ fechado sob subfórmulas e seja
  $\mModel{M} = \tuple{W, R, V}$. Uma \emph{filtração de $\mModel{M}$
  por $\Gamma$} é qualquer modelo
  $\mModel{M^*} = \tuple{W^*,R^*,V^*}$, onde:
  \begin{enumerate}
  \item $W^* = \Setabs{[w]}{w \in W}$;
  \item \ollabel{defn:filtration-R}%
    Para quaisquer $u,v \in W$:
    \begin{enumerate}
    \item \ollabel{defn:filtration-R1}%
      Se $Ruv$, então $R^*[u][v]$;
    \tagitem{prvBox}{\ollabel{defn:filtration-R2}%
      Se $R^*[u][v]$, então, para qualquer $\Box !A \in \Gamma$, se
      $\mSat{M}{\Box!A}[u]$, então $\mSat{M}{!A}[v]$}{};
    \tagitem{prvDiamond}{\ollabel{defn:filtration-R3}%
      Se $R^*[u][v]$, então, para qualquer $\Diamond !A \in \Gamma$, se
      $\mSat{M}{!A}[v]$, então $\mSat{M}{\Diamond!A}[u]$.}{}
    \end{enumerate}
  \item $V^*(p) = \Setabs{[u]}{u \in V(p)}$.
  \end{enumerate}
\end{defn}

Vale a pena pensar no que é $V^*(p)$: o conjunto constituído pelas
classes de equivalência~$[w]$ de todos os mundos~$w$ nos quais $p$ é
verdadeiro em~$\mModel{M}$. Por um lado, se $w \in V(p)$, então
$[w] \in V^*(p)$ por essa definição. Contudo, não é necessariamente
verdade que, se $[w] \in V^*(p)$, então $w \in V(p)$. Se
$[w] \in V^*(p)$, temos apenas a garantia de que $[w] = [u]$ para
\emph{algum}~$u \in V(p)$. É claro que $[w] = [u]$ significa que
$w \equiv u$. Portanto, quando $[w] \in V^*(p)$, podemos concluir
(apenas) que $w \equiv u$ para algum $u \in V(p)$.

\begin{thm}\ollabel{thm:filtrations}
  Se $\mModel{M^*}$ é uma filtração de $\mModel{M}$ por $\Gamma$,
  então, para todo $!A \in \Gamma$ e $w \in W$, temos
  $\mSat{M}{!A}[w]$ se e somente se $\mSat{M^*}{!A}[{[w]}]$.
\end{thm}

\begin{proof}
  Por indução em $!A$, usando o fato de que $\Gamma$ é fechado sob
  subfórmulas. Como $!A \in \Gamma$ e $\Gamma$ é fechado sob
  sub-!!{formula}s, todas as sub-!!{formula}s de~$!A$ também estão
  $\in \Gamma$. Logo, em cada passo indutivo, a hipótese de indução se
  aplica às sub-!!{formula}s de~$!A$.
  \begin{enumerate}
    \tagitem{prvFalse}{\indcase{!A}{\lfalse}{Não vale nem
        $\mSat{M}{\indfrm}[w]$ nem $\mSat{M^*}{\indfrm}[{[w]}]$.}}{}
    \tagitem{prvTrue}{\indcase{!A}{\ltrue}{Valem tanto
        $\mSat{M}{\indfrm}[w]$ quanto $\mSat{M^*}{\indfrm}[{[w]}]$.}}{}
    \item \indcase{!A}{p}{A direção da esquerda para a direita é
      imediata, pois $\mSat{M}{\indfrm}[w]$ somente se $w \in V(p)$,
      o que implica $[w] \in V^*(p)$, isto é,
      $\mSat{M^*}{\indfrm}[{[w]}]$. Reciprocamente, suponha
      $\mSat{M^*}{\indfrm}[{[w]}]$, isto é, $[w] \in V^*(p)$. Então,
      para algum $v \in V(p)$, $w \equiv v$. É claro que também
      $\mSat{M}{p}[v]$. Como $w \equiv v$, $w$ e $v$ tornam verdadeiras
      as mesmas !!{formula}s de~$\Gamma$. Como, por hipótese,
      $p \in \Gamma$ e $\mSat{M}{p}[v]$, $\mSat{M}{\indfrm}[w]$.}

    \tagitem{prvNot}{\indcase{!A}{\lnot !B}
         {\iftag{probNot}{Exercício.}
         {$\mSat{M}{\indfrm}[w]$ se e somente se
          $\mSat/{M}{!B}[w]$. Pela hipótese de indução,
          $\mSat/{M}{!B}[w]$ se e somente se
          $\mSat/{M^*}{!B}[{[w]}]$. Finalmente,
          $\mSat/{M^*}{!B}[{[w]}]$ se e somente se
          $\mSat{M^*}{\indfrm}[{[w]}]$.}}}{}
        
    \tagitem{prvAnd}{\indcase{!A}{(!B \land !C)}
          {\iftag{probAnd}{Exercício.}
          {$\mSat{M}{\indfrm}[w]$ se e somente se $\mSat{M}{!B}[w]$ e
          $\mSat{M}{!C}[w]$. Pela hipótese de indução, $\mSat{M}{!B}[w]$
          se e somente se $\mSat{M^*}{!B}[{[w]}]$, e $\mSat{M}{!C}[w]$ se e somente se
          $\mSat{M^*}{!C}[{[w]}]$. E $\mSat{M^*}{\indfrm}[{[w]}]$
          se e somente se $\mSat{M^*}{!B}[{[w]}]$ e
          $\mSat{M^*}{!C}[{[w]}]$.}}}{}

    \tagitem{prvOr}{\indcase{!A}{(!B \lor !C)}
          {\iftag{probOr}{Exercício.}
          {$\mSat{M}{\indfrm}[w]$ se e somente se $\mSat{M}{!B}[w]$ ou
          $\mSat{M}{!C}[w]$. Pela hipótese de indução, $\mSat{M}{!B}[w]$
          se e somente se $\mSat{M^*}{!B}[{[w]}]$, e $\mSat{M}{!C}[w]$ se e somente se
          $\mSat{M^*}{!C}[{[w]}]$. E $\mSat{M^*}{\indfrm}[{[w]}]$
          se e somente se $\mSat{M^*}{!B}[{[w]}]$ ou
          $\mSat{M^*}{!C}[{[w]}]$.}}}{}

    \tagitem{prvIf}{\indcase{!A}{(!B \lif !C)}
          {\iftag{probIf}{Exercício.}
          {$\mSat{M}{\indfrm}[w]$ se e somente se $\mSat/{M}{!B}[w]$ ou
          $\mSat{M}{!C}[w]$. Pela hipótese de indução, $\mSat{M}{!B}[w]$
          se e somente se $\mSat{M^*}{!B}[{[w]}]$, e $\mSat{M}{!C}[w]$ se e somente se
          $\mSat{M^*}{!C}[{[w]}]$. E $\mSat{M^*}{\indfrm}[{[w]}]$
          se e somente se $\mSat/{M^*}{!B}[{[w]}]$ ou
          $\mSat{M^*}{!C}[{[w]}]$.}}}{}

    \tagitem{prvIff}{\indcase{!A}{(!B \liff !C)}
          {\iftag{probIff}{Exercício.}
          {$\mSat{M}{\indfrm}[w]$ se e somente se $\mSat{M}{!B}[w]$ e
          $\mSat{M}{!C}[w]$, ou $\mSat/{M}{!B}[w]$ e
          $\mSat/{M}{!C}[w]$. Pela hipótese de indução,
          $\mSat{M}{!B}[w]$ se e somente se $\mSat{M^*}{!B}[{[w]}]$, e
          $\mSat{M}{!C}[w]$ se e somente se $\mSat{M^*}{!C}[{[w]}]$. E
          $\mSat{M^*}{\indfrm}[{[w]}]$ se e somente se
          $\mSat{M^*}{!B}[{[w]}]$ e $\mSat{M^*}{!C}[{[w]}]$, ou
          $\mSat/{M^*}{!B}[{[w]}]$ e
          $\mSat/{M^*}{!C}[{[w]}]$.}}}{}

    \tagitem{prvBox}{\indcase{!A}{\Box !B}
          {\iftag{probBox}{Exercício.}
          {Suponha $\mSat{M}{\indfrm}[w]$; para mostrar que
          $\mSat{M^*}{\indfrm}[{[w]}]$, seja $v$ tal que
          $R^*[w][v]$. Por
          \olref{defn:filtration}\olref{defn:filtration-R2}, temos
          $\mSat{M}{!B}[v]$ e, pela hipótese de indução,
          $\mSat{M^*}{!B}[{[v]}]$. Como $v$ era arbitrário, segue
          $\mSat{M^*}{\indfrm}[{[w]}]$.

          Reciprocamente, suponha $\mSat{M^*}{\indfrm}[{[w]}]$ e seja
          $v$ arbitrário tal que $Rwv$. Por
          \olref{defn:filtration}\olref{defn:filtration-R1}, temos
          $R^*[w][v]$, de modo que $\mSat{M^*}{!B}[{[v]}]$; pela
          hipótese de indução, $\mSat{M}{!B}[v]$ e, como $v$ era arbitrário,
          $\mSat{M}{\indfrm}[w]$.}}}{}
    
    \tagitem{prvDiamond}{\indcase{!A}{\Diamond !B}
          {\iftag{probDiamond}{Exercício.}
          {Suponha $\mSat{M}{\indfrm}[w]$. Então, para algum
          $v \in W$, $Rwv$ e $\mSat{M}{!B}[v]$. Pela hipótese de indução,
          $\mSat{M^*}{!B}[{[v]}]$ e, por
          \olref{defn:filtration}\olref{defn:filtration-R1}, temos
          $R^*[w][v]$. Assim, $\mSat{M^*}{\indfrm}[{[w]}]$.

          Suponha agora $\mSat{M^*}{\indfrm}[{[w]}]$. Então, para algum
          $[v] \in W^*$ com $R^*[w][v]$,
          $\mSat{M^*}{!B}[{[v]}]$. Pela hipótese de indução,
          $\mSat{M}{!B}[v]$. Por
          \olref{defn:filtration}\olref{defn:filtration-R3}, temos
          $\mSat{M}{\indfrm}[w]$.}}}{}
  \end{enumerate}
\end{proof}

\begin{prob}
  Complete a demonstração de \olref[nml][fil][fil]{thm:filtrations}
\end{prob}

O que vale para a verdade nos mundos de um modelo também vale para a
verdade em um modelo e para a validade em uma classe de modelos.

\begin{cor}
  Seja $\Gamma$ fechado sob subfórmulas. Então:
  \begin{enumerate}
  \item Se $\mModel{M^*}$ é uma filtração de $\mModel{M}$ por
    $\Gamma$, então, para qualquer $!A \in \Gamma$:
    $\mSat{M}{!A}$ se e somente se $\mSat{M^*}{!A}$.
  \item Se $\mClass{C}$ é uma classe de modelos e
    $\Gamma(\mClass{C})$ é a classe das $\Gamma$-filtrações de modelos
    em $\mClass{C}$, então qualquer !!{formula}~$!A \in \Gamma$ é
    válida em $\mClass{C}$ se e somente se é válida em
    $\Gamma(\mClass{C})$.
  \end{enumerate}
\end{cor}



\end{document}
